\section*{Exposición Aguda a Cuerpo Total}

\textbf{Dosis absorbida (rads)} \quad \textbf{Efecto}
\begin{itemize}
    \item \textbf{de 10} \quad Ruptura cromosómica difícil de detectar en células sanguíneas. Interferencia con la organogénesis de embriones.
    \item \textbf{de 25} \quad Cambios sanguíneos.
    \item \textbf{de 50} \quad Probable retención momentánea de la espermatogénesis.
    \item \textbf{de 100} \quad Síndrome de radiación probable.
    \item \textbf{de 200} \quad Síndrome grave de muerte.
    \item \textbf{de 400} \quad 50\% probabilidades de muerte.
    \item \textbf{de 600} \quad 100\% probabilidades de muerte.
\end{itemize}

Los síntomas del síndrome de radiación son: náuseas, vómitos, agotamiento, disminución de glóbulos blancos, perdida de recuperación aparente, vómito y diarrea con sangre, disminución de glóbulos blancos y plaquetas, anemia y finalmente intoxicación general.

Se llama exposición aguda cuando la dosis se recibe en un tiempo corto, por lo cual el cuerpo tiene menor oportunidad de reparación, y exposición crónica cuando la dosis se recibe en un tiempo largo.

El tiempo que tardan en presentarse los síntomas depende de la magnitud de la dosis y del intervalo de tiempo en que se imparta esta (dosis aguda o dosis crónica), es decir, cuando la dosis se imparte en un tiempo mayor el efecto es menor, como si el organismo recibiera una dosis menor.

En el caso de que la irradiación no sea en todo el cuerpo sino localizada en una parte de el, el daño depende de la magnitud del área y de los órganos expuestos, algunos de estos efectos, cuando la dosis se recibe en una sola exposición se describen a continuación.

\subsection*{Partes del Cuerpo Expuesto}
\begin{tabular}{lll}
\textbf{Parte del cuerpo expuesto} & \textbf{Dosis} & \textbf{Daño} \\
Gónadas & 50 rads & Esterilidad temporal\\
Gónadas & 800 rads & Esterilidad definitiva\\
Cuero cabelludo & 500 rads & Caída temporal del cabello\\
Cuero cabelludo & 2500 rads & Caída definitiva del cabello\\
\end{tabular}